\newlecture{1}{}{March 8, 2026}{Author Name}

Lecture 1: Introduction to the Course

\begin{definition}
  A \textbf{definition} is a statement that explains the meaning of a term or concept. It is used to clarify the meaning of a term and to provide a precise description of it.
\end{definition}

\begin{proposition}
  A \textbf{proposition} is a statement that can be proven to be true or false. It is a declarative sentence that asserts a fact or a relationship between concepts.
\end{proposition}

\begin{theorem}
  A \textbf{theorem} is a statement that has been proven to be true based on previously established statements, such as other theorems, axioms, and definitions. The proof of a theorem is a logical argument that demonstrates the truth of the statement.
\end{theorem}

\begin{lemma}
  A \textbf{lemma} is a preliminary result that is used to prove a theorem. It is a statement that is proven to be true and is often used as a stepping stone in the proof of a more complex theorem.
\end{lemma}

\begin{corollary}
  A \textbf{corollary} is a statement that follows directly from a theorem or a proposition. It is a consequence of the theorem and can be proven using the same reasoning as the original theorem.
\end{corollary}

\begin{remark}
  A \textbf{remark} is a comment or observation that is made about a theorem, proposition, definition, or any other mathematical statement. It is often used to provide additional insight, clarify a point, or highlight an interesting aspect of the statement.
\end{remark}

\begin{example}
  An \textbf{example} is a specific instance or case that illustrates a concept, definition, theorem, or any other mathematical statement. It is used to help understand the statement and to provide a concrete representation of it.
\end{example}

\begin{proof}
  A \textbf{proof} is a logical argument that demonstrates the truth of a mathematical statement. It is a sequence of statements that are logically connected and that lead to the conclusion that the statement is true. A proof can be written in various styles, such as direct proof, proof by contradiction, or proof by induction.
\end{proof}